\scanpage{6}
% archive sheet number: 5
% lines 1
$\bigl[$\uline{Lemme} A \quad $X$ espace annelé, \deleted{avec $\underline O_X$ \uline{cohérent}} \inserted{on suppose}.

% lines 2
(i) $X$ quasi-compact

% lines 3
\deleted{\unsure{(ii)} $\underline O_X$ cohérent} (ii) $\underline O_X$ cohérent, les $\underline O_{X,x}$ noethériens.

% lines 4–6
\deleted{\unsure{(ii)}} \inserted{(iii)} Pour \deleted{\unsure{tt} \unsure{suite exacte stricte}} \inserted{tout} Module \deleted{cohérent} \inserted{cohérent} $\underline J$ \deleted{\unsure{de} \unsure{dans} $\underline O_X$}, \inserted{sur $X$,} on a $H^1(X,J)=0$ \deleted{dans les cas suivants :} \deleted{$\underline J$ est libre de type fini \unsure{ou|sur les} \unsure{loc} \unsure{un}}

% lines 7
Alors, posant $\deleted{\unsure{N}}A=\Gamma(X,\underline O_X)$, le foncteur

% lines 8
\[
 M\unsure{\leadsto|\sim}\widetilde M=M\otimes_{A_X}\underline O_X
\]

% lines 9–12
de la catégorie des $A$-modules de présentation finie dans la catégorie des $\underline O_X$-Modules de prés. finie (= cohérents) est une équivalence de catégories. $\bigr]$

% lines left of 1–12
% margin note written sideways
\margin{\uline{Variante} Si tt pt de $X$ a un système fond. de voisinages \deleted{\unsure{qui}} quasi-compacts, si $X$ \deleted{étant} quasi-compact, et si pour tt \deleted{\unsure{sous}} Module $F$ \unsure{Noyau|d'Noyau} ou image d'un $\underline O^p\to\underline O^q$, $H^1(X,F)=0$, alors le foncteur $M\leadsto M\otimes_{A_X}\underline O_X$ des Modules sur $A$ dans les Modules \deleted{\unsure{quasi-coh}} ayant une \deleted{\unsure{une}} présentation \deleted{finie} \deleted{\unsure{globale}} \deleted{\unsure{est une}} \inserted{globale $\underline O^{(I)}-\underline O^{\unsure{(J)|J}}$, est} une équivalence \ldots}

% lines 13–15
\uline{Corollaire 1} Le foncteur $M\leadsto\widetilde M$ est \uline{exact}. En particulier, pour tt $x\in X$, $\underline O_{X,x}$ est \uline{plat} sur $A$.

% lines 16–18
\uline{Corollaire 2} L'anneau $A$ est \textquotedblleft cohérent\textquotedblright{} : i.e. pour tout hom $A^n - A$, le noyau est de t.f.

% lines 19–20
\uline{Corollaire 3} Si $X$ est \textquotedblleft pur\textquotedblright, \deleted{\unsure{les} $\underline O_{X,x}$} alors on a correspondance biunivoque entre idéaux maximaux de \uline{type fini} de $A$, et couples $(x,\mathfrak m)$ où $x\in X$ \deleted{\unsure{fermé}} et $\mathfrak m$ idéal maximal de $\underline O_{X,x}$.

% lines 21
\uline{Remarque} On peut remplacer (ii) (iii) par

% lines 22–23
(\deleted{\unsure{i}}$\alpha$) Pour toute suite exacte $\underline O_X^p-\underline O_X^q-F\to0$, on a exactitude pour la suite des $\Gamma$

% lines 24–25
(\deleted{\unsure{ii}}$\beta$) Tout $F$ de présentation finie est engendré par ses sections \ldots
